\part{Polyxan Hypergraphs and Media Quine Theory}

Part~II develops the algebraic and computational foundations of Polyxan
hypergraphs: typed, multi-layered, self-rewriting semantic structures
capable of representing every form of media on equal footing. Here the
theory of \emph{media quines} emerges: hypergraph structures that rewrite
themselves while preserving semantic invariants encoded in their types,
tags, and curvature signatures.

The chapters introduce typed morphisms, hyperedge signatures, quine
atoms, fibrations, curvature measures, spectral semantics, and sheaf-based
gluing conditions. Together, they define the structural conditions under
which Polyxan graphs can evolve, repair, replicate, and cohere under
semantic constraints.

These chapters also formalize the Polycompiler---the engine that extracts,
rewrites, hoists, reduces, compiles, and recombines Polyxan structures.
The Polycompiler is simultaneously a proof system, a media compiler, and
a semantic synthesizer.

Part~II concludes with the categorical foundations of Polyxan:
colimits, global polystructures, functorial compilers, and the limits of
semantic reconstruction, completing the algebraic half of the unified
RSVP--Polyxan theory.

\cleardoublepage
